\vskip 2cm

{\noindent\bfseries\Huge Resumen}

\vskip 1cm

Este documento presenta un Trabajo de Fin de Grado para la carrera de ingeniería informática en la Universidad de Valladolid.

El proyecto desarrolla una técnica de aprendizaje alternativa basada en el empleo de videojuegos con el fin de afianzar los conocimientos de manera dinámica y entretenida. Para ello se diseña un motor de videojuegos en Elixir; un videojuego didáctico, empleando también Elixir, haciendo uso del motor; una serie de prácticas diseñadas para alumnos de segundo año, que cursen la asignatura de Sistemas Distribuidos de la propia carrera mencionada al principio; posibles soluciones a la prácticas empleando Java y una página web con documentación sobre el desarrollo, haciendo uso del motor diseñado, de videojuegos para desarrolladores y diseñadores de videojuegos.

Todo el proyecto se diseña con la asignatura de Sistemas Distribuidos en mente, por lo tanto, el motor de videojuegos cuenta con la capacidad de ejecutarse de manera distribuida. Con el fin de ilustrar esta característica, el videojuego se diseña con partes distribuidas. Además se emplea Docker con el fin de simular un sistema distribuido dentro de una sola máquina.

\textbf{Palabras clave:} Sistema Distribuido, Elixir, Java, videojuego, aprendizaje.

\newpage

{\bfseries\Huge Abstract}

\vskip 1cm

This document presents a Bachelor's Final Project for the Computer Engineering degree at the University of Valladolid.

The project develops an alternative learning technique based on the use of videogames to consolidate knowledge in a dynamic and entertaining way. To this end, the following elements were designed: a videogame engine in Elixir; a didactic videogame, also using Elixir and utilizing the engine; a series of practical exercises designed for second-year students taking the Distributed Systems course of the aforementioned degree; possible solutions to the exercises using Java and a website containing documentation on videogame development using the designed engine for developers and designers.

The entire project is designed with the Distributed Systems course in mind; therefore, the videogame engine has the capability to run in a distributed manner. To illustrate this feature, the videogame is designed with distributed components. Furthermore, Docker is used to simulate a distributed system within a single machine.

\textbf{Keywords:} Distributed System, Elixir, Java, videogame, learning.